\section{一時所得}

\subsection{一時所得の計算}

\begin{fillblank}[年金1 2018 問題1(3)\textcircled{3}]
  確定給付企業年金の終了に伴い、その加入者であった A さんが平成30 年に残余財産の分配
金を一時金で受け取る場合の当該一時金にかかる課税所得は、（　　A　　）万円となる。

＜A さんの情報＞
\begin{itemize}
  \item 年齢：55 歳（在職者）
\item 確定給付企業年金の加入者期間：20 年
\item 残余財産の分配金額：1,080 万円
\item 確定給付企業年金において加入者負担掛金はない。
\item 分配金以外に課税対象となる所得は給与所得のみとし、税額計算における総所得に対する
各種所得控除（基礎控除等）は考慮しないものとする。
\end{itemize}

\end{fillblank}

\begin{answerbox}
  A:$(1,080万円 - 50万円) \times 1/2 = 515万円$
\end{answerbox}

\begin{explanationbox}
  DBの終了に係る残余財産の分配金は一時所得。
\end{explanationbox}

\newpage

\subsection{退職所得 or 一時所得}

\begin{truefalse}[年金1 2018 問題2(3)\textcircled{3}]
         \textbf{以下の文章の正誤を判定し、誤っている場合は正しい内容に改めなさい。}
\\  


  企業型年金の脱退一時金は、退職に起因する場合であっても一時所得となる。
\end{truefalse}

\begin{answerbox}
  \textbf{正}
\end{answerbox}

\begin{explanationbox}
  退職に起因して、老齢給付金を一時金で受給する場合は退職所得となる。
\end{explanationbox}

\begin{truefalse}[年金1 2019 問題2(1)\textcircled{1}]
         \textbf{以下の文章の正誤を判定し、誤っている場合は正しい内容に改めなさい。}
\\  


加入者拠出のある企業型年金において、一時金として支給される老齢給付金は、加入者拠出
分を除いて一時所得として課税される。
\end{truefalse}

\begin{answerbox}
  \textbf{誤}

  加入者拠出分も含めて退職所得として課税される。
\end{answerbox}

\begin{truefalse}[年金1 H29 問題1(6)\textcircled{4}]
         \textbf{下線部の正誤を判定し、誤っている場合は正しい内容に改めなさい。}
\\  

確定給付企業年金の老齢給付金を年金受給中の者が、その年金の一部を選択一時金として受
け取った場合、\underline{一時所得}として課税される。
\end{truefalse}

\begin{answerbox}
  \textbf{正}
\end{answerbox}

\begin{explanationbox}
  年金受給中に選択一時金を取得する場合は退職に起因しないので退職所得ではなく、一時所得。
\end{explanationbox}

\newpage